\chapter{开发工具、调试过程与实验验证}

\section{开发环境}

\begin{table}[htbp]
  \centering
  \caption{当前开发与验证工具}
  \begin{tabularx}{\textwidth}{p{3.2cm}p{4.1cm}X}
    \toprule
    工具 & 用途 & 定稿需记录 \\
    \midrule
    IAR Embedded Workbench & CYT4BB7双核工程编译与调试 & 版本、编译配置、优化等级 \\
    VOFA+ & JustFloat波形、在线调参和诊断 & 波特率115200，记录命令和通道 \\
    Python视觉上位机 & 使用真实固件C代码离线处理188×120帧 & Python/编译器版本、依赖锁定 \\
    摄像头标定工具 & 鱼眼标定、查表和固件头文件生成 & OpenCV版本、数据集与参数 \\
    RotorPy接口 & 对齐固件级联PID的动力学仿真 & 飞行器参数、噪声/延迟假设 \\
    IAR C-SPY & 在线下载与断点调试 & 调试器型号和下载流程 \\
    \bottomrule
  \end{tabularx}
\end{table}

\section{分层调试流程}

建议按照以下顺序记录，避免把上游观测问题误判为PID问题：
\begin{enumerate}
  \item 检查供电、电机方向、传感器轴向与串口连续性；
  \item 完成IMU六面安装矩阵和摄像头标定；
  \item 调整角速度环，再调姿态角环；
  \item 调整高度环、光流速度环；
  \item 地面验证视觉位置方向后，从低增益开始调位置环；
  \item 接入地面车模速度前馈和任务状态机；
  \item 在限高、限速、可急停条件下完成全链路联调。
\end{enumerate}

\section{VOFA可观测性设计}

当前工程提供从全局飞行总览、IMU、三轴角速度、姿态、高度、光流、位置、制动余量、视觉耗时到控制实时性的诊断模式。
报告不必枚举全部通道，但应选择能证明设计结论的波形，并保留原始CSV及采样频率。

\section{仿真与实机一致性}

RotorPy仿真复现位置—速度—姿态—角速度—混控链路。当前模型仍使用通用机体参数，
因此仿真只能验证控制语义和趋势；在用于量化结论前，应替换真实质量、轴距、惯量、推力系数、油门—推力曲线、观测噪声和延迟。

\section{核心实验与验收指标}

\begin{longtable}{p{3.0cm}p{4.0cm}p{3.5cm}p{2.4cm}}
\caption{整机实验矩阵}\label{tab:test-matrix}\\
\toprule
实验 & 自变量/工况 & 指标 & 数据状态 \\
\midrule
\endfirsthead
\toprule
实验 & 自变量/工况 & 指标 & 数据状态 \\
\midrule
\endhead
IMU静态与振动 & 静止、不同油门、悬停 & 噪声RMS、频谱峰值、姿态误差 & 待采集 \\
高度阶跃 & 不同目标高度与扰动 & 超调、调节时间、稳态误差 & 待采集 \\
水平速度跟踪 & X/Y/斜向阶跃 & RMSE、峰值、延迟、饱和 & 待采集 \\
位置跟踪 & 不同初始偏差与航向 & RMSE、停车误差、收敛时间 & 待采集 \\
视觉定位 & 光照、距离、姿态 & 检出率、误差、耗时 & 部分已有 \\
空地协同 & 正常、丢包、超时 & 完成时间、成功率、恢复时间 & 待采集 \\
整场任务 & 多次完整比赛流程 & 得分、完成时间、失败原因 & 待采集 \\
\bottomrule
\end{longtable}

\todoitem{每张曲线标明横纵轴、单位、采样率、测试条件和数据量；至少报告均值、最大值及重复次数，不只展示“最好的一次”。}
